\documentclass[twoside,final]{hcmut-report}


\coursename{Software Testing}
\reporttype{Project 2:}
\title{Black-box testing}
\advisor{& Bùi Hoài Thắng &}
\groupname{& 11 &}
\stuname{%
  & Lâm Minh Tùng Quân &  2352991 \\
  & Trương Xuân Sang &   2353045\\
  & Trương Gia Kỳ Nam &  2352787 \\
  & Nguyễn Hữu Phúc &  2352938 \\
  & Châu Anh Nhật  & 2352856\\
}


\begin{document}

\coverpage%

\tableofcontents
\pagebreak
\begin{table}[ht]
    \centering
    \caption{Project Contribution and Peer Review Mapping}
    \vspace{0.2cm}
    
    % Thiết lập độ rộng bảng bằng chiều rộng văn bản
    % Sử dụng cỡ chữ nhỏ hơn để vừa nội dung
    \small 
    \renewcommand{\arraystretch}{1.4} % Tăng chiều cao dòng để dễ đọc
    
    % Định nghĩa cột:
    % c: căn giữa (cho STT, %, ID)
    % p{...}: độ rộng cố định (cho Tên)
    % >{\raggedright\arraybackslash}X: Cột tự động co giãn, căn trái, cho phép xuống dòng
    
    \begin{tabularx}{\textwidth}{|c|c|p{2.5cm}|>{\raggedright\arraybackslash}X|c|}
    \hline
    \textbf{No.} & \textbf{ID} & \textbf{Full Name} & \textbf{Description} & \textbf{\%} \\ \hline
    
    1 & 2352938 & Nguyễn Hữu Phúc & 
 
    \begin{filelist}
        \item Feature 001: Admin Adds a New User 
    \end{filelist} & 20\% \\ \hline
    
    2 & 2352787 & Trương Gia Kỳ Nam & 

    \begin{filelist}
       \item Feature 006: Teacher Sets Up a Quiz
    \end{filelist} & 20\% \\ \hline
    
    3 & 2352856 & Châu Anh Nhật & 

    \begin{filelist}
        \item Feature 005: User Creates Calendar Event
    \end{filelist} & 20\% \\ \hline
    
    4 & 2353045 & Trương Xuân Sang & 

    \begin{filelist}
       \item Feature 002: Admin Creates a New Course
    \end{filelist} & 20\% \\ \hline
    
    5 & 2352991 & Lâm Minh Tùng Quân & 
    \begin{filelist}
       \item Feature 003: Teacher Creates Assignment
        \item Feature 004: Teacher Grades an Assignment 
    \end{filelist} & 20\% \\ \hline
    
    \end{tabularx}
\end{table}
\newpage
\newpage
% ==========================================
% 1. GENERAL INFORMATION
% ==========================================
\section{General Information}

\subsection{Selected Application}
Initially, the intended target for this testing campaign was the \textit{Mount Orange University} demonstration site. However, during the initial exploration phase, the default Moodle demo environment (version 5.2) exhibited critical system errors—specifically when attempting to add new "Activities or resources" to course sections. These server-side bugs severely impeded the execution of our core test scenarios.

In strict adherence to the project guidelines (Section: "Moodle demo goes wrong!"), which explicitly permits the registration of a new independent trial site to bypass these instability issues, our team migrated the testing environment. 

Consequently, the officially selected application for this report is a dedicated MoodleCloud LMS instance, hosted at: 
\begin{center}
    \textbf{\url{https://ihatetesting.moodlecloud.com/}}
\end{center}

\textit{Note: All test cases, data preparation, and automated executions detailed in this report were successfully performed on this stable environment, utilizing custom-created administrative and user accounts.}

\subsection{AI Tools Usage Declaration}
In strict compliance with the project guidelines, our team declares the utilization of \textbf{Google Gemini (Free Tier / Non-professional version)}. The usage and results of this tool are clearly defined as follows:

\begin{itemize}
    \item \textbf{Scope of Use:} The AI functioned strictly as a supportive assistant to brainstorm boundary values, structure complex \LaTeX{} tables, and refine the grammatical accuracy of the report.
    \item \textbf{Generated Results:} The direct outputs from the AI consisted of initial textual drafts and text-format diagrams (such as Mermaid code). 
    \item \textbf{Diagram Conversion:} All text-based diagrams generated by the AI were manually interpreted and rigorously redrawn into graphic formats using diagramming software.
    \item \textbf{Human Verification:} The core intellectual work—including testing logic design, equivalence class partitioning, test scenario formulation, and Katalon Recorder executions—was entirely manually performed and critically verified by human team members.
\end{itemize}

% ==========================================
% 2. REQUIREMENT SPECIFICATION
% ==========================================
\section{Requirement Specification}

\subsection{List of Selected Features}
\begin{itemize}
    \item \textbf{Feature 001:} Admin Adds a New User
    \item \textbf{Feature 002:} Admin Creates a New Course
    \item \textbf{Feature 003:} Teacher Creates an Assignment
    \item \textbf{Feature 004:} Teacher Grades an Assignment
    \item \textbf{Feature 005:} User Creates a Calendar Event
    \item \textbf{Feature 006:} Teacher Sets Up a Quiz
\end{itemize}

\subsection{Testing Environment}
To ensure the consistency, stability, and reproducibility of the test execution, all manual and automated testing campaigns were conducted under the following standardized hardware and software environment:

\begin{itemize}
    \item \textbf{Operating System:} Windows 11 Pro (64-bit, Version 23H2)
    \item \textbf{Web Browser:} Google Chrome (Version 124.0.6367.118 - Official Build)
    \item \textbf{Screen Resolution:} 1920 $\times$ 1080 pixels (Full HD) at 100\% display scale \textit{(Ensures consistent UI rendering during automated test execution)}.
    \item \textbf{Automation Tool:} Katalon Recorder (Chrome Extension, Version 5.9.0).
    \item \textbf{Network Condition:} Stable broadband internet connection \textit{(Crucial for avoiding false-negative timeout errors while interacting with the MoodleCloud server)}.
\end{itemize}

\textit{Note: The exact versions of the browser and extension were locked during the testing phase to prevent unexpected behavior caused by automatic updates.}

\newpage
% ==========================================
% 3. DETAILED TEST DESIGN
% ==========================================
\section{Detailed Test Design}

% --- FEATURE 001 ---
\subsection{Feature 001: Admin Adds a New User}

\subsubsection{Feature Description \& Usage Instruction}
The Site Administration module allows managers to manually create new user accounts in the Moodle system. This feature contains complex input validations regarding length, character formats, and password complexity policies.
\newline
\textbf{Usage Instruction:}
\begin{enumerate}
    \item Log in to \texttt{https://ihatetesting.moodlecloud.com/} \\ with username \texttt{phuc.nguyen0310@hcmut.edu.vn} and password \texttt{Huuphuc0310@}.
    \item Switch role to Manager.
    \item Navigate to \textbf{Site administration $>$ Users $>$ Accounts $>$ Add a new user}.
    \item Fill in the required and optional fields.
    \item Click \textbf{Create user}.
\end{enumerate}

\subsubsection{Data Preparation \& Clean-up (Setup/Teardown)}
To ensure test independence and prevent data clutter during automated test execution, the following data lifecycle scripts must be executed for this feature:
\begin{itemize}
    \item \textbf{Setup Script:} Log in to \texttt{https://ihatetesting.moodlecloud.com/}\\ with username \texttt{phuc.nguyen0310@hcmut.edu.vn} and password \texttt{Huuphuc0310@}. Switch role to Manager. Search the system to verify that the target test usernames do not already exist. If they do, delete them to prevent unintended duplicate errors during the "Happy Path" test execution.
    \item \textbf{Teardown Script:} Navigate to \textbf{Site administration $>$ Users $>$ Bulk user actions}. Click \textbf{Add all} to move the created test accounts into the "Users in list" section. Next, open the "With selected users..." dropdown, select \textbf{Delete}, and click \textbf{Go}. On the confirmation page, click \textbf{Yes} to permanently clear the data. Finally, \textbf{log out} of the system.
\end{itemize}

\subsubsection{Form Fields Catalog}
\begin{table}[H]
\centering
\renewcommand{\arraystretch}{1.3}
\begin{tabularx}{\textwidth}{|l|l|c|X|}
\hline
\textbf{Field} & \textbf{Type} & \textbf{Req.} & \textbf{Constraints / Business Rules} \\ \hline
Username & Text input & Yes & Lowercase letters, digits, hyphens, underscores, periods, @. No spaces. Must be unique. \\ \hline
Authentication method & Dropdown & Yes & Default: ``Manual accounts''. Options: Manual, Email-based self-registration, etc. \\ \hline
Suspended account & Checkbox & No & Default: unchecked. \\ \hline
Generate password & Checkbox & No & If checked, manual password input is bypassed and system notifies the user. \\ \hline
New password & Text input & Cond. & Min 8 chars, $\geq$1 digit, $\geq$1 lowercase, $\geq$1 uppercase, $\geq$1 non-alphanumeric char. Max 128 chars. \\ \hline
Force password change & Checkbox & No & Default: unchecked. If checked, user must reset password on first login. \\ \hline
First name & Text input & Yes & Cannot be empty. Max 100 chars. \\ \hline
Last name & Text input & Yes & Cannot be empty. Max 100 chars. \\ \hline
Email address & Text input & Yes & Must be valid email format. Max 100 chars. Duplicates rejected unless allowed. \\ \hline
\end{tabularx}
\end{table}

\subsubsection{Boundary Value Analysis (BVA)}
\textbf{Methodology Applied: Robust Boundary Value Analysis (Robust BVA).} 
Based on the single-fault assumption theory, for a variable bounded by $[min, max]$, Robust BVA normally generates exactly 7 test values. \textbf{Critical Execution Note:} When the target variable is tested at a boundary, all other variables in the form are explicitly held at their valid nominal ($nom$) values to prevent fault masking.

\textbf{Variables Analyzed Summary:}
\begin{itemize}
    \item \textbf{V1: Username Length} (4 values: $min-, min, min+, nom$ - UI lacks maximum limit).
    \item \textbf{V2: Password Length} (6 values: $min-, min, min+, max-, max, max+$).
    \item \textbf{V3: First Name Length} (5 values: $min-, min, min+, max-, max$ - UI caps at 100, $max+$ impossible).
    \item \textbf{V4: Last Name Length} (5 values: $min-, min, min+, max-, max$ - UI caps at 100, $max+$ impossible).
\end{itemize}

\vspace{0.3cm}
\noindent\textbf{Variable 1: Username Length} \\
\textit{Domain: $1 \le \text{Length}$}
\begin{table}[H]
\centering
\renewcommand{\arraystretch}{1.2}
\begin{tabularx}{\textwidth}{|c|c|l|X|}
\hline
\textbf{Boundary} & \textbf{Length} & \textbf{Target Variable Test Input} & \textbf{Expected Result} \\ \hline
$min-$ & 0 & \textit{(empty string)} & ❌ Error: Form submission fails \\ \hline
$min$ & 1 & \texttt{usr003\_u} & ✅ Accepted: User created successfully \\ \hline
$min+$ & 2 & \texttt{usr004\_uu} & ✅ Accepted: User created successfully \\ \hline
$nom$ & 17 & \texttt{usr001\_uuuuuuuuuu} & ✅ Accepted: User created successfully \\ \hline
\end{tabularx}
\end{table}

\vspace{0.3cm}
\noindent\textbf{Variable 2: Password Length} \\
\textit{Domain: $8 \le \text{Length} \le 128$ (Moodle site policy, verified by UI testing)}
\begin{table}[H]
\centering
\renewcommand{\arraystretch}{1.4}
\begin{tabularx}{\textwidth}{|c|c|X|X|}
\hline
\textbf{Boundary} & \textbf{Len} & \textbf{Target Variable Test Input} & \textbf{Expected Result} \\ \hline
$min-$ & 7 & \texttt{\small Ab1!cD2} & ❌ Error: Passwords must be at least 8 chars \\ \hline
$min$ & 8 & \texttt{\small Ab1!cD2@} & ✅ Accepted: User created successfully \\ \hline
$min+$ & 9 & \texttt{\small Ab1!cD2@a} & ✅ Accepted: User created successfully \\ \hline
$max-$ & 127 & \texttt{\tiny Ab1!cD2@eF3\#gH4\allowbreak{}\$iJ5\%kL6\textasciicircum{}mN7\&oP\allowbreak{}8*qR9=sT0+uVwXy\allowbreak{}ZAb1!cD2@eF3\#gH\allowbreak{}4\$iJ5\%kL6\textasciicircum{}mN7\&o\allowbreak{}P8*qR9=sT0+uVwX\allowbreak{}yZAb1!cD2@eF3\#g\allowbreak{}H4\$iJ5\%kL6\textasciicircum{}mN7\&\allowbreak{}oP8*qR} & ✅ Accepted: User created successfully \\ \hline
$max$ & 128 & \texttt{\tiny Ab1!cD2@eF3\#gH4\allowbreak{}\$iJ5\%kL6\textasciicircum{}mN7\&oP\allowbreak{}8*qR9=sT0+uVwXy\allowbreak{}ZAb1!cD2@eF3\#gH\allowbreak{}4\$iJ5\%kL6\textasciicircum{}mN7\&o\allowbreak{}P8*qR9=sT0+uVwX\allowbreak{}yZAb1!cD2@eF3\#g\allowbreak{}H4\$iJ5\%kL6\textasciicircum{}mN7\&\allowbreak{}oP8*qR9=} & ✅ Accepted: User created successfully \\ \hline
$max+$ & 129 & \texttt{\tiny Ab1!cD2@eF3\#gH4\allowbreak{}\$iJ5\%kL6\textasciicircum{}mN7\&oP\allowbreak{}8*qR9=sT0+uVwXy\allowbreak{}ZAb1!cD2@eF3\#gH\allowbreak{}4\$iJ5\%kL6\textasciicircum{}mN7\&o\allowbreak{}P8*qR9=sT0+uVwX\allowbreak{}yZAb1!cD2@eF3\#g\allowbreak{}H4\$iJ5\%kL6\textasciicircum{}mN7\&\allowbreak{}oP8*qR9=s} & ❌ Error: Password exceeds maximum length \\ \hline
\end{tabularx}
\end{table}

\vspace{0.3cm}
\noindent\textbf{Variable 3 \& 4: First Name \& Last Name Length} \\
\textit{Domain: $1 \le \text{Length} \le 100$}
\begin{table}[H]
\centering
\renewcommand{\arraystretch}{1.2}
\begin{tabularx}{\textwidth}{|c|c|l|X|}
\hline
\textbf{Boundary} & \textbf{Length} & \textbf{Target Variable Test Input} & \textbf{Expected Result} \\ \hline
$min-$ & 0 & \textit{(empty string)} & ❌ Error: Form submission fails \\ \hline
$min$ & 1 & \texttt{f} / \texttt{l} & ✅ Accepted: User created successfully \\ \hline
$min+$ & 2 & \texttt{ff} / \texttt{ll} & ✅ Accepted: User created successfully \\ \hline
$max-$ & 99 & \texttt{f...} / \texttt{l...} (99 chars) & ✅ Accepted: User created successfully \\ \hline
$max$ & 100 & \texttt{f...} / \texttt{l...} (100 chars) & ✅ Accepted: User created successfully \\ \hline
\end{tabularx}
\end{table}

\subsubsection{Equivalence Class Partitioning (ECP)}
\textbf{Methodology Applied: Weak Robust ECP.}
The input domain is partitioned into mutually exclusive Valid and Invalid classes. One representative value is selected from each class. Similar to BVA, the Single-Fault Assumption is strictly enforced during execution.

\textbf{Variables Analyzed Summary:}
\begin{itemize}
    \item \textbf{V1: Username Character Format} (4 Invalid Classes).
    \item \textbf{V2: Password Character Format} (4 Invalid Classes).
    \item \textbf{V3: Email Address Format} (3 Invalid Classes).
\end{itemize}

\vspace{0.3cm}
\noindent\textbf{Equivalence Classes Definition:}
\begin{table}[H]
\centering
\renewcommand{\arraystretch}{1.2}
\begin{tabularx}{\textwidth}{|c|l|X|l|X|}
\hline
\textbf{Var} & \textbf{ID} & \textbf{Equivalence Class} & \textbf{Representative} & \textbf{Expected Result} \\ \hline
\multirow{4}{*}{\textbf{User}} 
& U5 & Contains uppercase & \texttt{usr026\_UsrNom} & ❌ Error near Username \\ \cline{2-5}
& U6 & Contains spaces & \texttt{usr027\_usr nom} & ❌ Error near Username \\ \cline{2-5}
& U7 & Special characters & \texttt{usr028\_usr!nom} & ❌ Error near Username \\ \cline{2-5}
& U9 & Existing duplicate & \texttt{admin} & ❌ Error near Username \\ \hline
\multirow{4}{*}{\textbf{Pass}} 
& P2 & Missing numeric digit & \texttt{Abcdefg!@} & ❌ Error near Password \\ \cline{2-5}
& P3 & Missing uppercase & \texttt{abcdefg1!@} & ❌ Error near Password \\ \cline{2-5}
& P4 & Missing lowercase & \texttt{ABCDEFG1!@} & ❌ Error near Password \\ \cline{2-5}
& P5 & Missing special char & \texttt{Abcdefg123} & ❌ Error near Password \\ \hline
\multirow{3}{*}{\textbf{Email}} 
& E4 & Missing @ symbol & \texttt{testexample.com} & ❌ Error near Email \\ \cline{2-5}
& E5 & Missing domain & \texttt{test@} & ❌ Error near Email \\ \cline{2-5}
& E7 & Contains spaces & \texttt{te st@ex.com} & ❌ Error near Email \\ \hline
\end{tabularx}
\end{table}

\subsubsection{Decision Table Testing}
\textbf{Methodology Applied: Algebraic Simplification \& Rule Counting.} 
The Password Complexity Policy relies on 5 Boolean conditions ($C_1$ to $C_5$). A complete Limited Entry Decision Table (LEDT) would have $2^5 = 32$ rules. To optimize testing without losing coverage, we apply algebraic simplification using "Don't Care" entries (\texttt{-}). 

\begin{table}[H]
\centering
\renewcommand{\arraystretch}{1.2}
\begin{tabularx}{\textwidth}{|>{\raggedright\arraybackslash}X|c|c|c|c|c|c|}
\hline
\textbf{Conditions} & \textbf{R1} & \textbf{R2} & \textbf{R3} & \textbf{R4} & \textbf{R5} & \textbf{R6} \\ \hline
C1: Length $\geq$ 8 & T & T & T & T & T & F \\ \hline
C2: Has $\geq$1 digit & T & T & T & T & F & - \\ \hline
C3: Has $\geq$1 uppercase & T & T & T & F & - & - \\ \hline
C4: Has $\geq$1 lowercase & T & T & F & - & - & - \\ \hline
C5: Has $\geq$1 special char & T & F & - & - & - & - \\ \hline
\hline
\textbf{Actions / System Output} & & & & & & \\ \hline
A1: Accept password & \textbf{X} & & & & & \\ \hline
A2: Reject \& show error message & & \textbf{X} & \textbf{X} & \textbf{X} & \textbf{X} & \textbf{X} \\ \hline
\end{tabularx}
\end{table}

\subsubsection{Use-Case Testing}
\textbf{Methodology Applied: Path-based Scenario Generation.} 
Test scenarios are systematically derived by traversing distinct execution paths (Basic, Alternative, and Exception paths) identified in the Activity Diagram.

\vspace{0.2cm}
\noindent\textbf{Use Case Name:} Admin Adds a New User \\
\textbf{Actor:} Site Administrator (Manager) \\
\textbf{Pre-condition:} 
\begin{enumerate}
    \item Log in to \texttt{https://ihatetesting.moodlecloud.com/} \\ with username \texttt{phuc.nguyen0310@hcmut.edu.vn} and password \texttt{Huuphuc0310@}.
    \item Switch role to Manager.
    \item Navigate to: \textbf{Site administration $>$ Users $>$ Accounts $>$ Add a new user}.
\end{enumerate}
\textbf{Post-condition:} A new user record is inserted into the Moodle database, followed by the user logging out of the system.

\vspace{0.3cm}
\noindent\textbf{Derived Test Scenarios:}
\begin{table}[H]
\centering
\renewcommand{\arraystretch}{1.2}
\begin{tabularx}{\textwidth}{|l|l|X|}
\hline
\textbf{Scenario ID} & \textbf{Path Type} & \textbf{Description \& Expected Result} \\ \hline
S1 & Basic Flow (Happy Path) & Admin enters all valid data and a manual password. \textbf{Result:} User is created successfully. \\ \hline
S2 & Alternative Flow & Admin checks "Generate password and notify user". \textbf{Result:} Auto-password generated, user created. \\ \hline
S3 & Exception Flow & Duplicate username (U9). \textbf{Result:} Form submission fails, error shown near Username. \\ \hline
S4 & Exception Flow & Invalid password. \textbf{Result:} Form submission fails, error shown near Password. \\ \hline
S5 & Exception Flow & Invalid email. \textbf{Result:} Form submission fails, error shown near Email address. \\ \hline
S6 & Exception Flow & Empty required field. \textbf{Result:} Form submission fails, error shown near Required fields. \\ \hline
\end{tabularx}
\end{table}

\begin{figure}[H]
    \centering
    % Chèn hình ảnh graphics.TC-001.png vào đây
    \includegraphics[width=0.8\textwidth]{graphics/TC-001.png}
    \caption{UML Activity Diagram outlining the execution paths for User Creation}
    \label{fig:user-creation-activity}
\end{figure}

\subsubsection{Test Cases Summary (Excerpt)}
\textit{Note: This is an excerpt of key test cases. The complete set of test cases, detailing multi-variable data combinations, is documented in the attached \texttt{Testcases.xlsx} file.}

\begin{table}[H]
\centering
\renewcommand{\arraystretch}{1.2}
\begin{tabularx}{\textwidth}{|l|c|X|X|}
\hline
\textbf{TC ID} & \textbf{Method} & \textbf{Test Case Name / Coverage} & \textbf{Expected Result} \\ \hline
TC-001-002 & BVA & V1: Username (min-) & ❌ Error: Form submission fails \\ \hline
TC-001-016 & BVA & V4: Last Name (min-) & ❌ Error: Form submission fails \\ \hline
TC-001-021 & ECP & Username (U5: Uppercase) & ❌ Error: Form submission fails \\ \hline
TC-001-025 & ECP & Password (P2: No digit) & ❌ Error: Form submission fails \\ \hline
TC-001-039 & DT  & R2: Missing non-alphanumeric character & ❌ Error: Form submission fails \\ \hline
\multicolumn{4}{|c|}{\cellcolor{gray!20}\textbf{Use-Case Scenarios Coverage}} \\ \hline
TC-001-032 & UC & Cover S1: Happy path & ✅ Accepted: User created successfully \\ \hline
TC-001-033 & UC & Cover S2: Generate password & ✅ Accepted: User created successfully \\ \hline
TC-001-034 & UC & Cover S3: Duplicate username & ❌ Error: Form submission fails \\ \hline
TC-001-037 & UC & Cover S6: Empty required field & ❌ Error: Form submission fails \\ \hline
\end{tabularx}
\end{table}

% --- FEATURE 002 ---
\subsection{Feature 002: Admin Creates a New Course}

\subsubsection{Feature Description \& Usage Instruction}
The Course Management module allows managers or administrators to create new course spaces with configurable settings. This feature enforces chronological constraints on dates, limits on text lengths, and uniqueness rules for course short names.
\newline
\textbf{Usage Instruction:}
\begin{enumerate}
    \item Log in to \texttt{https://ihatetesting.moodlecloud.com/}\\ with username \texttt{phuc.nguyen0310@hcmut.edu.vn} and password \texttt{Huuphuc0310@}.
    \item Switch role to Manager.
    \item Navigate to \textbf{Site administration $>$ Courses $>$ Add a new course}.
    \item Fill in the Course full name, Short name, dates, format, and other settings.
    \item Click \textbf{Save and display}.
\end{enumerate}

\subsubsection{Data Preparation \& Clean-up (Setup/Teardown)}
To ensure test independence and prevent data clutter during automated test execution, the following data lifecycle scripts must be executed for this feature:
\begin{itemize}
    \item \textbf{Setup Script:} Log in to \texttt{https://ihatetesting.moodlecloud.com/}\\ with username \texttt{phuc.nguyen0310@hcmut.edu.vn} and password \texttt{Huuphuc0310@}. Switch role to Manager. Search the system to verify that the target test course short names (e.g., \texttt{sn001\_ssssssss}) do not already exist. If they do, delete them to prevent unintended duplicate errors during test execution.
    \item \textbf{Teardown Script:} Navigate to \textbf{Site administration $>$ Courses $>$ Manage courses and categories}. Locate the category where the test courses were created. Select the test courses, click the \textbf{Delete} (trash bin) icon, and click \textbf{Continue} on the confirmation page to permanently clear the data. Finally, \textbf{log out} of the system.
\end{itemize}

\subsubsection{Form Fields Catalog}
\begin{table}[H]
\centering
\renewcommand{\arraystretch}{1.3}
\begin{tabularx}{\textwidth}{|l|l|c|X|}
\hline
\textbf{Field} & \textbf{Type} & \textbf{Req.} & \textbf{Constraints / Business Rules} \\ \hline
Course full name & Text input & Yes & Cannot be empty. Max 254 chars. \\ \hline
Course short name & Text input & Yes & Must be unique across all courses. Max 255 chars. \\ \hline
Course category & Dropdown & Yes & Must select a valid existing category. \\ \hline
Course start date & Date picker & Yes & Must be a valid date. \\ \hline
Enable end date & Checkbox & No & Toggles end date validation logic. \\ \hline
Course end date & Date picker & Cond. & Must be strictly AFTER the start date if "Enable" is checked. \\ \hline
Course format & Dropdown & Yes & Options: Weekly, Topics, Social, Single activity. \\ \hline
\end{tabularx}
\end{table}

\subsubsection{Boundary Value Analysis (BVA)}
\textbf{Methodology Applied: Robust Boundary Value Analysis (Robust BVA).} 
Based on the single-fault assumption theory, for a variable bounded by $[min, max]$, Robust BVA normally generates exactly 7 test values. \textbf{Critical Execution Note:} When the target variable is tested at a boundary, all other variables in the form are explicitly held at their valid nominal ($nom$) values to prevent fault masking.

\textbf{Variables Analyzed Summary:}
\begin{itemize}
    \item \textbf{V1: Full Name Length} (5 values: $min-, min, min+, max-, max$ - UI caps input, $max+$ impossible).
    \item \textbf{V2: Short Name Length} (5 values: $min-, min, min+, max-, max$ - UI caps input, $max+$ impossible).
    \item \textbf{V3: End Date Chronology} (6 values: $min-, min, min+, max-, max, max+$ relative to Start Date).
\end{itemize}

\vspace{0.3cm}
\noindent\textbf{Variable 1: Full Name Length} \\
\textit{Domain: $1 \le \text{Length} \le 254$ (UI max; $max+$ impossible and omitted)}
\begin{table}[H]
\centering
\renewcommand{\arraystretch}{1.4}
\begin{tabularx}{\textwidth}{|c|c|X|X|}
\hline
\textbf{Boundary} & \textbf{Len} & \textbf{Target Variable Test Input} & \textbf{Expected Result} \\ \hline
$min-$ & 0 & \textit{(empty string)} & ❌ Error: Form submission fails \\ \hline
$min$ & 1 & \texttt{\small f} & ✅ Accepted: Course created successfully \\ \hline
$min+$ & 2 & \texttt{\small fn} & ✅ Accepted: Course created successfully \\ \hline
$max-$ & 253 & \texttt{\tiny fn005\_nnnnnnnnn\allowbreak{}nnnnnnnnnnnnnnn\allowbreak{}nnnnnnnnnnnnnnn\allowbreak{}nnnnnnnnnnnnnnn\allowbreak{}nnnnnnnnnnnnnnn\allowbreak{}nnnnnnnnnnnnnnn\allowbreak{}nnnnnnnnnnnnnnn\allowbreak{}nnnnnnnnnnnnnnn\allowbreak{}nnnnnnnnnnnnnnn\allowbreak{}nnnnnnnnnnnnnnn\allowbreak{}nnnnnnnnnnnnnnn\allowbreak{}nnnnnnnnnnnnnnn\allowbreak{}nnnnnnnnnnnnnnn\allowbreak{}nnnnnnnnnnnnnnn\allowbreak{}nnnnnnnnnnnnnnn\allowbreak{}nnnnnnnnnnnnnnn\allowbreak{}nnnnnnnnnnnnn} & ✅ Accepted: Course created successfully \\ \hline
$max$ & 254 & \texttt{\tiny fn006\_nnnnnnnnn\allowbreak{}nnnnnnnnnnnnnnn\allowbreak{}nnnnnnnnnnnnnnn\allowbreak{}nnnnnnnnnnnnnnn\allowbreak{}nnnnnnnnnnnnnnn\allowbreak{}nnnnnnnnnnnnnnn\allowbreak{}nnnnnnnnnnnnnnn\allowbreak{}nnnnnnnnnnnnnnn\allowbreak{}nnnnnnnnnnnnnnn\allowbreak{}nnnnnnnnnnnnnnn\allowbreak{}nnnnnnnnnnnnnnn\allowbreak{}nnnnnnnnnnnnnnn\allowbreak{}nnnnnnnnnnnnnnn\allowbreak{}nnnnnnnnnnnnnnn\allowbreak{}nnnnnnnnnnnnnnn\allowbreak{}nnnnnnnnnnnnnnn\allowbreak{}nnnnnnnnnnnnnn} & ✅ Accepted: Course created successfully \\ \hline
\end{tabularx}
\end{table}

\vspace{0.3cm}
\noindent\textbf{Variable 2: Short Name Length} \\
\textit{Domain: $1 \le \text{Length} \le 255$ (UI max; $max+$ impossible and omitted)}
\begin{table}[H]
\centering
\renewcommand{\arraystretch}{1.4}
\begin{tabularx}{\textwidth}{|c|c|X|X|}
\hline
\textbf{Boundary} & \textbf{Len} & \textbf{Target Variable Test Input} & \textbf{Expected Result} \\ \hline
$min-$ & 0 & \textit{(empty string)} & ❌ Error: Form submission fails \\ \hline
$min$ & 1 & \texttt{\small s} & ✅ Accepted: Course created successfully \\ \hline
$min+$ & 2 & \texttt{\small sn} & ✅ Accepted: Course created successfully \\ \hline
$max-$ & 254 & \texttt{\tiny sn011\_sssssssss\allowbreak{}sssssssssssssss\allowbreak{}sssssssssssssss\allowbreak{}sssssssssssssss\allowbreak{}sssssssssssssss\allowbreak{}sssssssssssssss\allowbreak{}sssssssssssssss\allowbreak{}sssssssssssssss\allowbreak{}sssssssssssssss\allowbreak{}sssssssssssssss\allowbreak{}sssssssssssssss\allowbreak{}sssssssssssssss\allowbreak{}sssssssssssssss\allowbreak{}sssssssssssssss\allowbreak{}sssssssssssssss\allowbreak{}sssssssssssssss\allowbreak{}ssssssssssssss} & ✅ Accepted: Course created successfully \\ \hline
$max$ & 255 & \texttt{\tiny sn012\_sssssssss\allowbreak{}sssssssssssssss\allowbreak{}sssssssssssssss\allowbreak{}sssssssssssssss\allowbreak{}sssssssssssssss\allowbreak{}sssssssssssssss\allowbreak{}sssssssssssssss\allowbreak{}sssssssssssssss\allowbreak{}sssssssssssssss\allowbreak{}sssssssssssssss\allowbreak{}sssssssssssssss\allowbreak{}sssssssssssssss\allowbreak{}sssssssssssssss\allowbreak{}sssssssssssssss\allowbreak{}sssssssssssssss\allowbreak{}sssssssssssssss\allowbreak{}sssssssssssssss\allowbreak{}} & ✅ Accepted: Course created successfully \\ \hline
\end{tabularx}
\end{table}

\vspace{0.3cm}
\noindent\textbf{Variable 3: End Date Chronology (Relative to Start Date)} \\
\textit{Domain: $\text{End Date} > \text{Start Date}$}
\begin{table}[H]
\centering
\renewcommand{\arraystretch}{1.2}
\begin{tabularx}{\textwidth}{|c|l|X|}
\hline
\textbf{Boundary} & \textbf{Target Variable Test Input} & \textbf{Expected Result} \\ \hline
$min-$ & \texttt{Start Date - 1 day} & ❌ Error: End date must be after start date \\ \hline
$min$ & \texttt{Exact Start Date} & ❌ Error: End date must be after start date \\ \hline
$min+$ & \texttt{Start Date + 1 day} & ✅ Accepted: Course created successfully \\ \hline
$max-$ & \texttt{Start Date + 10 years} & ✅ Accepted: Course created successfully \\ \hline
$max$ & \texttt{Start Date + 11 years} & ✅ Accepted: Course created successfully \\ \hline
$max+$ & \texttt{Start Date + 12 years} & ✅ Accepted: Course created successfully \\ \hline
\end{tabularx}
\end{table}

\subsubsection{Equivalence Class Partitioning (ECP)}
\textbf{Methodology Applied: Weak Robust ECP.}
The input domain is partitioned into mutually exclusive Valid and Invalid classes. One representative value is selected from each class. Similar to BVA, the Single-Fault Assumption is strictly enforced during execution.

\textbf{Variables Analyzed Summary:}
\begin{itemize}
    \item \textbf{V1: Short Name Uniqueness} (1 Invalid Class).
\end{itemize}

\vspace{0.3cm}
\noindent\textbf{Equivalence Classes Definition:}
\begin{table}[H]
\centering
\renewcommand{\arraystretch}{1.2}
\begin{tabularx}{\textwidth}{|c|l|X|l|X|}
\hline
\textbf{Var} & \textbf{ID} & \textbf{Equivalence Class} & \textbf{Representative} & \textbf{Expected Result} \\ \hline
\multirow{1}{*}{\textbf{S.Name}} 
& SN2 & Existing duplicate name & \texttt{TC001} (existing short name) & ❌ Error near Short name \\ \hline
\end{tabularx}
\end{table}

\subsubsection{Decision Table Testing}
\textbf{Methodology Applied: Algebraic Simplification \& Rule Counting.} 
The UI logic for End Date configuration is dependent on the selected Course Format and the Checkboxes. According to system documentation, "Calculate the end date from the number of sections" is available \textbf{for courses in weekly format only}.

\begin{table}[H]
\centering
\renewcommand{\arraystretch}{1.2}
\begin{tabularx}{\textwidth}{|>{\raggedright\arraybackslash}X|c|c|c|c|c|}
\hline
\textbf{Conditions} & \textbf{R1} & \textbf{R2} & \textbf{R3} & \textbf{R4} & \textbf{R5} \\ \hline
C1: Course format = Weekly & F & F & T & T & T \\ \hline
C2: End date Enable checked & F & T & F & T & T \\ \hline
C3: Calculate end date from sections checked & - & - & - & F & T \\ \hline
\hline
\textbf{Actions / System Output} & & & & & \\ \hline
A1: End date pickers enabled (editable) & & \textbf{X} & & \textbf{X} & \\ \hline
A2: End date pickers disabled (auto-calculated) & & & & & \textbf{X} \\ \hline
A3: End date pickers disabled (hidden/ignored) & \textbf{X} & & \textbf{X} & & \\ \hline
\end{tabularx}
\end{table}

\subsubsection{Use-Case Testing}
\textbf{Methodology Applied: Path-based Scenario Generation.} 
Test scenarios are systematically derived by traversing distinct execution paths (Basic, Alternative, and Exception paths) identified in the Activity Diagram.

\vspace{0.2cm}
\noindent\textbf{Use Case Name:} Admin Creates a New Course \\
\textbf{Actor:} Site Administrator (Manager) \\
\textbf{Pre-condition:} 
\begin{enumerate}
    \item Log in to \texttt{https://ihatetesting.moodlecloud.com/}\\ with username \texttt{phuc.nguyen0310@hcmut.edu.vn} and password \texttt{Huuphuc0310@}.
    \item Switch role to Manager.
    \item Navigate to: \textbf{Site administration $>$ Courses $>$ Add a new course}.
\end{enumerate}
\textbf{Post-condition:} A new course is created and visible in the Moodle database, followed by the user logging out of the system.

\vspace{0.3cm}
\noindent\textbf{Derived Test Scenarios:}
\begin{table}[H]
\centering
\renewcommand{\arraystretch}{1.2}
\begin{tabularx}{\textwidth}{|l|l|X|}
\hline
\textbf{Scenario ID} & \textbf{Path Type} & \textbf{Description \& Expected Result} \\ \hline
S1 & Basic Flow (Happy Path) & Admin enters valid data, enables end date (after start date). \textbf{Result:} Course is created successfully. \\ \hline
S2 & Alternative Flow (AF1) & Admin disables end date validation. \textbf{Result:} Course created successfully without end date limits. \\ \hline
S3 & Alternative Flow (AF2) & Admin sets course visibility to ``Hide. \textbf{Result:} Course is created but hidden from students. \\ \hline
S4 & Exception Flow (EF1) & Duplicate short name entered. \textbf{Result:} Form submission fails, error shown near Short name. \\ \hline
S5 & Exception Flow (EF2) & End date is set before start date (Date Conflict). \textbf{Result:} Form submission fails, error shown. \\ \hline
S6 & Exception Flow (EF3) & Empty required field (Full name or short name). \textbf{Result:} Form submission fails, error shown. \\ \hline
\end{tabularx}
\end{table}

\begin{figure}[H]
    \centering
    % Chèn hình ảnh graphics.TC-002.png vào đây
    \includegraphics[width=0.8\textwidth]{graphics/TC-002.png}
    \caption{UML Activity Diagram outlining the execution paths for Course Creation}
    \label{fig:course-creation-activity}
\end{figure}

\subsubsection{Test Cases Summary (Excerpt)}
\textit{Note: This is an excerpt of key test cases. The complete set of test cases, detailing multi-variable data combinations, is documented in the attached \texttt{Testcases.xlsx} file.}

\begin{table}[H]
\centering
\renewcommand{\arraystretch}{1.2}
\begin{tabularx}{\textwidth}{|l|c|X|X|}
\hline
\textbf{TC ID} & \textbf{Method} & \textbf{Test Case Name / Coverage} & \textbf{Expected Result} \\ \hline
TC-002-002 & BVA & V1: Full Name Length (min-) & ❌ Error: Form submission fails \\ \hline
TC-002-007 & BVA & V2: Short Name Length (min-) & ❌ Error: Form submission fails \\ \hline
TC-002-012 & BVA & V3: End Date Chronology (min-) & ❌ Error: Form submission fails \\ \hline
TC-002-018 & ECP & Short Name (SN2: Duplicate) & ❌ Error: Form submission fails \\ \hline
TC-002-023 & DT  & R5: Weekly + End Date ON + Auto-calc & ✅ Accepted: Course created successfully \\ \hline
\multicolumn{4}{|c|}{\cellcolor{gray!20}\textbf{Use-Case Scenarios Coverage}} \\ \hline
TC-002-024 & UC & Cover S1: Happy path & ✅ Accepted: Course created successfully \\ \hline
TC-002-025 & UC & Cover S2: End date disabled & ✅ Accepted: Course created successfully \\ \hline
TC-002-026 & UC & Cover S4: Duplicate short name & ❌ Error: Form submission fails \\ \hline
TC-002-027 & UC & Cover S5: Date conflict & ❌ Error: Form submission fails \\ \hline
\end{tabularx}
\end{table}

% --- FEATURE 003 ---
\subsection{Feature 003: Teacher Creates Assignment}

\subsubsection{Feature Description \& Usage Instruction}
The Assignment activity module enables teachers to collect work from students, review it, and provide grades and feedback. This feature contains complex validations regarding submission types, chronological date constraints (Allow submissions, Due date, Cut-off date), and grading limits.
\newline
\textbf{Usage Instruction:}
\begin{enumerate}
    \item Log in to \texttt{https://ihatetesting.moodlecloud.com/}\\ with username \texttt{phuc.nguyen0310@hcmut.edu.vn} and password \texttt{Huuphuc0310@}.
    \item Navigate to a specific course and toggle \textbf{Edit mode} on.
    \item Click \textbf{Add an activity or resource $>$ Assignment}.
    \item Configure the assignment name, dates, submission types, and grading settings.
    \item Click \textbf{Save and return to course}.
\end{enumerate}

\subsubsection{Data Preparation \& Clean-up (Setup/Teardown)}
To ensure test independence and prevent data clutter during automated test execution, the following data lifecycle scripts must be executed for this feature:
\begin{itemize}
    \item \textbf{Setup Script:} Log in to \texttt{https://ihatetesting.moodlecloud.com/} as a Teacher. Access the target test course. Ensure that no duplicate assignment names (e.g., \texttt{an001\_aaaaaaaa...}) exist in the course sections. If they do, delete them to maintain a clean test environment.
    \item \textbf{Teardown Script:} Navigate back to the course page and ensure \textbf{Edit mode} is toggled on. Locate the newly created test assignments. Click the 3-dots (More) menu next to each assignment, select \textbf{Delete}, and click \textbf{Yes} on the confirmation modal to permanently clear the data. Finally, \textbf{log out} of the system.
\end{itemize}

\subsubsection{Form Fields Catalog}
\begin{table}[H]
\centering
\renewcommand{\arraystretch}{1.3}
\begin{tabularx}{\textwidth}{|l|l|c|X|}
\hline
\textbf{Field} & \textbf{Type} & \textbf{Req.} & \textbf{Constraints / Business Rules} \\ \hline
Assignment name & Text input & Yes & Cannot be empty. Maximum length enforced by UI. \\ \hline
Allow submissions from & Date picker & No & Valid date format. \\ \hline
Due date & Date picker & No & Must be ON or AFTER the "Allow submissions from" date. \\ \hline
Cut-off date & Date picker & No & Must be ON or AFTER the "Due date". \\ \hline
File submissions & Checkbox & No & Enables/disables file upload capabilities. \\ \hline
Online text & Checkbox & No & Enables/disables inline text submissions. \\ \hline
Grade to pass & Number input & No & Value must be $\le$ Maximum grade (Default: 100). \\ \hline
\end{tabularx}
\end{table}

\subsubsection{Boundary Value Analysis (BVA)}
\textbf{Methodology Applied: Robust Boundary Value Analysis (Robust BVA).} 
Based on the single-fault assumption theory, for a variable bounded by $[min, max]$, Robust BVA normally generates exactly 7 test values. \textbf{Critical Execution Note:} When the target variable is tested at a boundary, all other variables in the form are explicitly held at their valid nominal ($nom$) values to prevent fault masking.

\textbf{Variables Analyzed Summary:}
\begin{itemize}
    \item \textbf{V1: Assignment Name Length} (5 values: $min-, min, min+, max-, max$ - UI caps input, $max+$ impossible).
    \item \textbf{V2: Grade to Pass} (3 values: $max-, max, max+$ - $min$ boundary has no lower limit UI constraint besides zero).
    \item \textbf{V3: Due Date Chronology} (5 values: $min-, min, min+, max-, max$ relative to Allow Date).
    \item \textbf{V4: Cut-off Date Chronology} (5 values: $min-, min, min+, max-, max$ relative to Due Date).
\end{itemize}

\vspace{0.3cm}
\noindent\textbf{Variable 1: Assignment Name Length} \\
\textit{Note: Assignment Name has a hard UI maximum limit, therefore the $max+$ boundary is impossible and omitted.}
\begin{table}[H]
\centering
\renewcommand{\arraystretch}{1.2}
\begin{tabularx}{\textwidth}{|c|c|X|X|}
\hline
\textbf{Boundary} & \textbf{Len} & \textbf{Target Variable Test Input} & \textbf{Expected Result} \\ \hline
$min-$ & 0 & \textit{(empty string)} & ❌ Error: Form submission fails \\ \hline
$min$ & 1 & \texttt{\small a} & ✅ Accepted: Assignment created \\ \hline
$min+$ & 2 & \texttt{\small as} & ✅ Accepted: Assignment created \\ \hline
$max-$ & 256 & \texttt{\tiny an005\_aaaaaaaaa\allowbreak{}aaaaaaaaaaaaaaa\allowbreak{}aaaaaaaaaaaaaaa\allowbreak{}aaaaaaaaaaaaaaa\allowbreak{}aaaaaaaaaaaaaaa\allowbreak{}aaaaaaaaaaaaaaa\allowbreak{}aaaaaaaaaaaaaaa\allowbreak{}aaaaaaaaaaaaaaa\allowbreak{}aaaaaaaaaaaaaaa\allowbreak{}aaaaaaaaaaaaaaa\allowbreak{}aaaaaaaaaaaaaaa\allowbreak{}aaaaaaaaaaaaaaa\allowbreak{}aaaaaaaaaaaaaaa\allowbreak{}aaaaaaaaaaaaaaa\allowbreak{}aaaaaaaaaaaaaaa\allowbreak{}aaaaaaaaaaaaaaa\allowbreak{}aaaaaaaaaaaaaaa\allowbreak{}a} & ✅ Accepted: Assignment created \\ \hline
$max$ & 257 & \texttt{\tiny an006\_aaaaaaaaa\allowbreak{}aaaaaaaaaaaaaaa\allowbreak{}aaaaaaaaaaaaaaa\allowbreak{}aaaaaaaaaaaaaaa\allowbreak{}aaaaaaaaaaaaaaa\allowbreak{}aaaaaaaaaaaaaaa\allowbreak{}aaaaaaaaaaaaaaa\allowbreak{}aaaaaaaaaaaaaaa\allowbreak{}aaaaaaaaaaaaaaa\allowbreak{}aaaaaaaaaaaaaaa\allowbreak{}aaaaaaaaaaaaaaa\allowbreak{}aaaaaaaaaaaaaaa\allowbreak{}aaaaaaaaaaaaaaa\allowbreak{}aaaaaaaaaaaaaaa\allowbreak{}aaaaaaaaaaaaaaa\allowbreak{}aaaaaaaaaaaaaaa\allowbreak{}aaaaaaaaaaaaaaa\allowbreak{}aa} & ✅ Accepted: Assignment created \\ \hline
\end{tabularx}
\end{table}

\vspace{0.3cm}
\noindent\textbf{Variable 2: Grade to Pass} \\
\textit{Domain: $\text{Grade} \le \text{Maximum Grade (100)}$}
\begin{table}[H]
\centering
\renewcommand{\arraystretch}{1.2}
\begin{tabularx}{\textwidth}{|c|l|X|}
\hline
\textbf{Boundary} & \textbf{Target Variable Test Input} & \textbf{Expected Result} \\ \hline
$max-$ & \texttt{99} & ✅ Accepted: Assignment created \\ \hline
$max$ & \texttt{100} & ✅ Accepted: Assignment created \\ \hline
$max+$ & \texttt{101} & ❌ Error: Grade to pass cannot be greater than maximum grade \\ \hline
\end{tabularx}
\end{table}

\vspace{0.3cm}
\noindent\textbf{Variable 3: Due Date Chronology (Relative to Allow Submissions Date)}
\begin{table}[H]
\centering
\renewcommand{\arraystretch}{1.2}
\begin{tabularx}{\textwidth}{|c|l|X|}
\hline
\textbf{Boundary} & \textbf{Target Variable Test Input} & \textbf{Expected Result} \\ \hline
$min-$ & \texttt{Allow Submissions Date (same day)} & ❌ Error: Due date cannot equal or precede Allow date \\ \hline
$min$ & \texttt{Allow Submissions Date + 1 day} & ✅ Accepted: Assignment created \\ \hline
$min+$ & \texttt{Allow Submissions Date + 2 days} & ✅ Accepted: Assignment created \\ \hline
$max-$ & \texttt{Allow Submissions Date + 10 years} & ✅ Accepted: Assignment created \\ \hline
$max$ & \texttt{Allow Submissions Date + 11 years} & ✅ Accepted: Assignment created \\ \hline
\end{tabularx}
\end{table}

\vspace{0.3cm}
\noindent\textbf{Variable 4: Cut-off Date Chronology (Relative to Due Date)}
\begin{table}[H]
\centering
\renewcommand{\arraystretch}{1.2}
\begin{tabularx}{\textwidth}{|c|l|X|}
\hline
\textbf{Boundary} & \textbf{Target Variable Test Input} & \textbf{Expected Result} \\ \hline
$min-$ & \texttt{Due Date - 1 day} & ❌ Error: Cut-off date cannot be earlier than Due date \\ \hline
$min$ & \texttt{Exact Due Date} & ✅ Accepted: Assignment created \\ \hline
$min+$ & \texttt{Due Date + 1 day} & ✅ Accepted: Assignment created \\ \hline
$max-$ & \texttt{Due Date + 10 years} & ✅ Accepted: Assignment created \\ \hline
$max$ & \texttt{Due Date + 11 years} & ✅ Accepted: Assignment created \\ \hline
\end{tabularx}
\end{table}

\subsubsection{Equivalence Class Partitioning (ECP)}
\textbf{Methodology Applied: Weak Robust ECP.}
The input domain is partitioned into mutually exclusive Valid and Invalid classes. For this specific feature, independent invalid ECP classes separate from BVA limits are non-existent. Constraints such as non-numeric inputs in grade fields or date format validation are implicitly handled by the UI (dropdowns/number pickers) and fully covered by the BVA boundaries detailed above.

\subsubsection{Decision Table Testing}
\textbf{Methodology Applied: Component Logic Mapping.} 
Submission settings depend on the combination of checkboxes selected by the teacher. These checkboxes control the visibility and validation of subsequent form fields dynamically on the UI.

\begin{table}[H]
\centering
\renewcommand{\arraystretch}{1.2}
\begin{tabularx}{\textwidth}{|>{\raggedright\arraybackslash}X|c|c|c|c|}
\hline
\textbf{Conditions} & \textbf{R1} & \textbf{R2} & \textbf{R3} & \textbf{R4} \\ \hline
C1: File submissions checked & T & T & F & F \\ \hline
C2: Online text checked & T & F & T & F \\ \hline
\hline
\textbf{Actions / System Output (UI Visibility)} & & & & \\ \hline
A1: Max files \& Submission size dropdowns visible & \textbf{X} & \textbf{X} & & \\ \hline
A2: Accepted file types field visible & \textbf{X} & \textbf{X} & & \\ \hline
A3: Word limit configuration visible & \textbf{X} & & \textbf{X} & \\ \hline
A4: Warning: No submission method selected & & & & \textbf{X} \\ \hline
\end{tabularx}
\end{table}

\subsubsection{Use-Case Testing}
\textbf{Methodology Applied: Path-based Scenario Generation.} 
Test scenarios are systematically derived by traversing distinct execution paths (Basic, Alternative, and Exception paths) identified in the Activity Diagram.

\vspace{0.2cm}
\noindent\textbf{Use Case Name:} Teacher Creates Assignment \\
\textbf{Actor:} Teacher \\
\textbf{Pre-condition:} 
\begin{enumerate}
    \item Log in to \texttt{https://ihatetesting.moodlecloud.com/}\\ with username \texttt{phuc.nguyen0310@hcmut.edu.vn} and password \texttt{Huuphuc0310@}.
    \item Navigate to the target course and toggle \textbf{Edit mode} on.
    \item Click \textbf{Add an activity or resource $>$ Assignment}.
\end{enumerate}
\textbf{Post-condition:} A new assignment activity is published within the course, followed by the user logging out of the system.

\vspace{0.3cm}
\noindent\textbf{Derived Test Scenarios:}
\begin{table}[H]
\centering
\renewcommand{\arraystretch}{1.2}
\begin{tabularx}{\textwidth}{|l|l|X|}
\hline
\textbf{Scenario ID} & \textbf{Path Type} & \textbf{Description \& Expected Result} \\ \hline
S1 & Basic Flow (Happy Path) & File submission enabled, valid due date set, Grade to Pass = 50. \textbf{Result:} Assignment created successfully. \\ \hline
S2 & Alternative Flow & Online text submission only, Due date disabled (open-ended). \textbf{Result:} Assignment created successfully. \\ \hline
S3 & Alternative Flow & Both Online text and File submissions enabled. \textbf{Result:} Assignment created successfully. \\ \hline
S4 & Exception Flow & Empty assignment name. \textbf{Result:} Form submission fails, error shown. \\ \hline
S5 & Exception Flow & Cut-off date set before due date. \textbf{Result:} Form submission fails, error shown. \\ \hline
S6 & Exception Flow & Grade to pass exceeds Maximum grade limit. \textbf{Result:} Form submission fails, error shown. \\ \hline
\end{tabularx}
\end{table}

\begin{figure}[H]
    \centering
    % Chèn hình ảnh graphics.TC-001.png vào đây
    \includegraphics[width=0.5\textwidth]{graphics/TC-003.png}
    \caption{UML Activity Diagram outlining the execution paths for Assignment Creation}
    \label{fig:assignment-creation-activity}
\end{figure}

\subsubsection{Test Cases Summary (Excerpt)}
\textit{Note: This is an excerpt of key test cases. The complete set of test cases, detailing multi-variable data combinations, is documented in the attached \texttt{Testcases.xlsx} file.}

\begin{table}[H]
\centering
\renewcommand{\arraystretch}{1.2}
\begin{tabularx}{\textwidth}{|l|c|X|X|}
\hline
\textbf{TC ID} & \textbf{Method} & \textbf{Test Case Name / Coverage} & \textbf{Expected Result} \\ \hline
TC-003-002 & BVA & V1: Name Length (min-) & ❌ Error: Form submission fails \\ \hline
TC-003-009 & BVA & V2: Grade to Pass (max+) & ❌ Error: Form submission fails \\ \hline
TC-003-010 & BVA & V3: Due vs Allow (min-) & ❌ Error: Form submission fails \\ \hline
TC-003-015 & BVA & V4: Cut-off vs Due (min-) & ❌ Error: Form submission fails \\ \hline
TC-003-021 & DT  & R2: File Submissions Only & ✅ Accepted: Assignment created \\ \hline
\multicolumn{4}{|c|}{\cellcolor{gray!20}\textbf{Use-Case Scenarios Coverage}} \\ \hline
TC-003-024 & UC & Cover S1: Happy path & ✅ Accepted: Assignment created \\ \hline
TC-003-025 & UC & Cover S2: No due date (Open-ended) & ✅ Accepted: Assignment created \\ \hline
TC-003-026 & UC & Cover S4: Empty name error & ❌ Error: Form submission fails \\ \hline
TC-003-027 & UC & Cover S3/S5: Date conflict & ❌ Error: Form submission fails \\ \hline
\end{tabularx}
\end{table}

% --- FEATURE 004 ---
\subsection{Feature 004: Teacher Grades an Assignment}

\subsubsection{Feature Description \& Usage Instruction}
The grading interface allows teachers to assign numerical scores and provide feedback on student submissions. This feature ensures that grades fall within the defined maximum grade limit (default out of 100) and triggers notifications to students.
\newline
\textbf{Usage Instruction:}
\begin{enumerate}
    \item Log in to \texttt{https://ihatetesting.moodlecloud.com/}\\ with username \texttt{phuc.nguyen0310@hcmut.edu.vn} and password \texttt{Huuphuc0310@}.
    \item Navigate to the target course and click on the target Assignment activity.
    \item Click \textbf{View all submissions}.
    \item Click the \textbf{Grade} button for a specific student.
    \item Enter a numeric value in the \textbf{Grade} field and/or add text in \textbf{Feedback comments}.
    \item Click \textbf{Save changes}.
\end{enumerate}

\subsubsection{Data Preparation \& Clean-up (Setup/Teardown)}
To ensure test independence and prevent data clutter during automated test execution, the following data lifecycle scripts must be executed for this feature:
\begin{itemize}
    \item \textbf{Setup Script:} Log in as a Teacher. Navigate to the target Assignment's grading page. Ensure that the target test student (e.g., \texttt{student01}) has a blank grade (not graded). If a grade already exists, clear the grade field and save before running the test cases to avoid state interference.
    \item \textbf{Teardown Script:} After the test execution finishes, remain on the grading page for the test student. Delete the value in the \textbf{Grade} field, clear any text in the \textbf{Feedback comments} editor, and click \textbf{Save changes} to revert the student's submission back to the "Not graded" state. Finally, \textbf{log out} of the system.
\end{itemize}

\subsubsection{Form Fields Catalog}
\begin{table}[H]
\centering
\renewcommand{\arraystretch}{1.3}
\begin{tabularx}{\textwidth}{|l|l|c|X|}
\hline
\textbf{Field} & \textbf{Type} & \textbf{Req.} & \textbf{Constraints / Business Rules} \\ \hline
Grade out of 100 & Number input & Yes* & Must be a numeric value between 0 and the maximum grade (100). *Can be empty if only leaving feedback. \\ \hline
Feedback comments & Rich Text & No & Allows formatted text, images, and links. \\ \hline
Notify student & Checkbox & No & If checked, an email/system notification is sent to the student upon saving. Default is checked. \\ \hline
\end{tabularx}
\end{table}

\subsubsection{Boundary Value Analysis (BVA)}
\textbf{Methodology Applied: Robust Boundary Value Analysis (Robust BVA).} 
Based on the single-fault assumption theory, for a variable bounded by $[min, max]$, Robust BVA generates exactly 7 test values. \textbf{Critical Execution Note:} When the target variable is tested at a boundary, all other variables in the form are explicitly held at their valid nominal ($nom$) values to prevent fault masking.

\textbf{Variables Analyzed Summary:}
\begin{itemize}
    \item \textbf{V1: Grade Value} (7 values: $min-, min, min+, nom, max-, max, max+$ - bounding domain $[0, 100]$).
\end{itemize}

\vspace{0.3cm}
\noindent\textbf{Variable 1: Grade Value} \\
\textit{Domain: $0 \le \text{Grade} \le 100$}
\begin{table}[H]
\centering
\renewcommand{\arraystretch}{1.2}
\begin{tabularx}{\textwidth}{|c|c|l|X|}
\hline
\textbf{Boundary} & \textbf{Value} & \textbf{Target Variable Test Input} & \textbf{Expected Result} \\ \hline
$min-$ & -1 & \texttt{-1} & ❌ Error: Grade must be 0 or more \\ \hline
$min$ & 0 & \texttt{0} & ✅ Accepted: Grade saved successfully \\ \hline
$min+$ & 1 & \texttt{1} & ✅ Accepted: Grade saved successfully \\ \hline
$nom$ & 50 & \texttt{50} & ✅ Accepted: Grade saved successfully \\ \hline
$max-$ & 99 & \texttt{99} & ✅ Accepted: Grade saved successfully \\ \hline
$max$ & 100 & \texttt{100} & ✅ Accepted: Grade saved successfully \\ \hline
$max+$ & 101 & \texttt{101} & ❌ Error: Grade must be 100 or less \\ \hline
\end{tabularx}
\end{table}

\subsubsection{Equivalence Class Partitioning (ECP)}
\textbf{Methodology Applied: Weak Robust ECP.}
The input domain is partitioned into mutually exclusive Valid and Invalid classes. One representative value is selected from each class. Similar to BVA, the Single-Fault Assumption is strictly enforced during execution.

\textbf{Variables Analyzed Summary:}
\begin{itemize}
    \item \textbf{V1: Grade Input Format} (5 Invalid Classes: G5–G9).
\end{itemize}

\vspace{0.3cm}
\noindent\textbf{Equivalence Classes Definition:}
\begin{table}[H]
\centering
\renewcommand{\arraystretch}{1.2}
\begin{tabularx}{\textwidth}{|c|l|X|l|X|}
\hline
\textbf{Var} & \textbf{ID} & \textbf{Equivalence Class} & \textbf{Representative} & \textbf{Expected Result} \\ \hline
\multirow{5}{*}{\textbf{Grade}}
& G5 & Negative value & \texttt{-5} & ❌ Error near Grade field \\ \cline{2-5}
& G6 & Above maximum ($>$100) & \texttt{105} & ❌ Error near Grade field \\ \cline{2-5}
& G7 & Alphabetic characters & \texttt{abc} & ❌ Error near Grade field \\ \cline{2-5}
& G8 & Special characters & \texttt{!@\#} & ❌ Error near Grade field \\ \cline{2-5}
& G9 & Empty (no grade entered) & \textit{(empty)} & ✅ Accepted: Feedback saved \\ \hline
\end{tabularx}
\end{table}

\subsubsection{Decision Table Testing}
\textbf{Methodology Applied: System Workflow Logic.} 
The grading workflow incorporates two primary conditions: the validity of the grade input and the state of the "Notify student" checkbox. If the grade input is invalid, the system rejects the submission entirely, making the notification state irrelevant ("Don't Care" rule applied).

\begin{table}[H]
\centering
\renewcommand{\arraystretch}{1.2}
\begin{tabularx}{\textwidth}{|>{\raggedright\arraybackslash}X|c|c|c|}
\hline
\textbf{Conditions} & \textbf{R1} & \textbf{R2} & \textbf{R3} \\ \hline
C1: Grade input is valid (0 $\le$ Value $\le$ 100) & T & T & F \\ \hline
C2: Notify student checkbox is checked & T & F & - \\ \hline
\hline
\textbf{Actions / System Output} & & & \\ \hline
A1: Save grade to database successfully & \textbf{X} & \textbf{X} & \\ \hline
A2: Trigger notification event to student & \textbf{X} & & \\ \hline
A3: Block save \& show validation error & & & \textbf{X} \\ \hline
\end{tabularx}
\end{table}

\subsubsection{Use-Case Testing}
\textbf{Methodology Applied: Path-based Scenario Generation.} 
Test scenarios are systematically derived by traversing distinct execution paths (Basic, Alternative, and Exception paths) identified in the Activity Diagram.

\vspace{0.2cm}
\noindent\textbf{Use Case Name:} Teacher Grades an Assignment \\
\textbf{Actor:} Teacher \\
\textbf{Pre-condition:} 
\begin{enumerate}
    \item Log in to \texttt{https://ihatetesting.moodlecloud.com/}\\ with username \texttt{phuc.nguyen0310@hcmut.edu.vn} and password \texttt{Huuphuc0310@}.
    \item Navigate to the target course $>$ Assignment activity.
    \item Click \textbf{View all submissions} and open the grading interface for a student.
\end{enumerate}
\textbf{Post-condition:} The student's grade and feedback are updated in the gradebook, followed by the user logging out of the system.

\vspace{0.3cm}
\noindent\textbf{Derived Test Scenarios:}
\begin{table}[H]
\centering
\renewcommand{\arraystretch}{1.2}
\begin{tabularx}{\textwidth}{|l|l|X|}
\hline
\textbf{Scenario ID} & \textbf{Path Type} & \textbf{Description \& Expected Result} \\ \hline
S1 & Basic Flow (Happy Path) & Teacher enters a valid numeric grade (50), adds feedback, leaves notify checked, and clicks \textbf{Save changes}. \textbf{Result:} Grade saved and student notified. \\ \hline
S2 & Alternative Flow (AF1) & Teacher clicks \textbf{Save and show next} to grade and advance to the next student. \textbf{Result:} Grade saved and grading panel advances. \\ \hline
S3 & Alternative Flow (AF2) & Teacher provides only feedback (Grade field empty), unchecks notify, and saves. \textbf{Result:} Feedback saved without grade, no notification sent. \\ \hline
S4 & Exception Flow (EF1) & Teacher enters a grade above 100 ($max+$). \textbf{Result:} Save blocked, error shown. \\ \hline
S5 & Exception Flow (EF2) & Teacher enters a grade below 0 ($min-$). \textbf{Result:} Save blocked, error shown. \\ \hline
S6 & Exception Flow (EF3) & Teacher enters non-numeric text in the grade field. \textbf{Result:} Save blocked, error shown. \\ \hline
\end{tabularx}
\end{table}

\begin{figure}[H]
    \centering
    % Chèn hình ảnh graphics.TC-004.png vào đây
    \includegraphics[width=0.8\textwidth]{graphics/TC-004.png}
    \caption{UML Activity Diagram outlining the execution paths for Assignment Grading}
    \label{fig:assignment-grading-activity}
\end{figure}

\subsubsection{Test Cases Summary (Excerpt)}
\textit{Note: This is an excerpt of key test cases. The complete set of test cases, detailing multi-variable data combinations, is documented in the attached \texttt{Testcases.xlsx} file.}

\begin{table}[H]
\centering
\renewcommand{\arraystretch}{1.2}
\begin{tabularx}{\textwidth}{|l|c|X|X|}
\hline
\textbf{TC ID} & \textbf{Method} & \textbf{Test Case Name / Coverage} & \textbf{Expected Result} \\ \hline
TC-004-002 & BVA & V1: Grade Value (min-) & ❌ Error: Form submission fails \\ \hline
TC-004-006 & BVA & V1: Grade Value (max) & ✅ Accepted: Grade saved successfully \\ \hline
TC-004-008 & ECP & Grade (G5: Negative value) & ❌ Error: Form submission fails \\ \hline
TC-004-010 & ECP & Grade (G7: Alphabetic) & ❌ Error: Form submission fails \\ \hline
TC-004-012 & ECP & Grade (G9: Empty — feedback only) & ✅ Accepted: Grade saved successfully \\ \hline
\multicolumn{4}{|c|}{\cellcolor{gray!20}\textbf{Use-Case Scenarios Coverage}} \\ \hline
TC-004-013 & UC & Cover S1: Happy path & ✅ Accepted: Grade saved successfully \\ \hline
TC-004-014 & UC & Cover S2: Save and show next & ✅ Accepted: Grade saved successfully \\ \hline
TC-004-015 & UC & Cover S3: Feedback only, no notify & ✅ Accepted: Grade saved successfully \\ \hline
TC-004-016 & UC & Cover S4: Grade over maximum & ❌ Error: Form submission fails \\ \hline
TC-004-017 & UC & Cover S6: Non-numeric grade & ❌ Error: Form submission fails \\ \hline
\end{tabularx}
\end{table}

% --- FEATURE 005 ---
\subsection{Feature 005: User Creates Calendar Event}

\subsubsection{Feature Description \& Usage Instruction}
The Calendar module allows users to schedule personal, group, or course-wide events. This feature enforces validation on event titles, date chronologies, and duration settings (either a specific end date or a time span in minutes).
\newline
\textbf{Usage Instruction:}
\begin{enumerate}
    \item Log in to \texttt{https://ihatetesting.moodlecloud.com/}\\ with username \texttt{phuc.nguyen0310@hcmut.edu.vn} and password \texttt{Huuphuc0310@}.
    \item Navigate to the \textbf{Calendar} page from the left navigation menu.
    \item Click the \textbf{New event} button.
    \item Fill in the Event title, Date, and configure the Duration settings.
    \item Click \textbf{Save}.
\end{enumerate}

\subsubsection{Data Preparation \& Clean-up (Setup/Teardown)}
To ensure test independence and prevent data clutter during automated test execution, the following data lifecycle scripts must be executed for this feature:
\begin{itemize}
    \item \textbf{Setup Script:} Log in with the designated credentials. Navigate to the Calendar. Verify that no test events with target names (e.g., \texttt{t001\_tttttttttttttttttttt}) exist on the current testing date. If they do, click on them and delete them prior to execution.
    \item \textbf{Teardown Script:} Locate the newly created test event on the calendar grid. Click the event block to open the details modal. Click the \textbf{Delete} button, then click \textbf{Delete event} on the confirmation prompt to permanently clear the data. Finally, \textbf{log out} of the system.
\end{itemize}

\subsubsection{Form Fields Catalog}
\begin{table}[H]
\centering
\renewcommand{\arraystretch}{1.3}
\begin{tabularx}{\textwidth}{|l|l|c|X|}
\hline
\textbf{Field} & \textbf{Type} & \textbf{Req.} & \textbf{Constraints / Business Rules} \\ \hline
Event title & Text input & Yes & Cannot be empty. Max length limits apply. \\ \hline
Date & Date/Time picker & Yes & Must be a valid calendar date and time. \\ \hline
Duration & Radio group & Yes & Options: Without duration, Until, Duration in minutes. Default: Without duration. \\ \hline
Until (End Date) & Date/Time picker & Cond. & Required if "Until" is selected. Must be ON or AFTER the start date. \\ \hline
Duration in minutes & Number input & Cond. & Required if "Duration in minutes" is selected. Must be $\geq 1$. \\ \hline
\end{tabularx}
\end{table}

\subsubsection{Boundary Value Analysis (BVA)}
\textbf{Methodology Applied: Robust Boundary Value Analysis (Robust BVA).} 
Based on the single-fault assumption theory, for a variable bounded by $[min, max]$, Robust BVA generates exactly 7 test values. \textbf{Critical Execution Note:} When the target variable is tested at a boundary, all other variables in the form are explicitly held at their valid nominal ($nom$) values to prevent fault masking.

\textbf{Variables Analyzed Summary:}
\begin{itemize}
    \item \textbf{V1: Event Title Length} (3 values: $min-, min, min+$ — no UI-enforced maximum; $max-$, $max$, $max+$ are impossible and omitted).
    \item \textbf{V2: Duration in Minutes} (7 values: $min-, min, min+, nom, max-, max, max+$ — UI enforces maximum of 9{,}999{,}999 minutes).
\end{itemize}

\vspace{0.3cm}
\noindent\textbf{Variable 1: Event Title Length} \\
\textit{Note: Event Title has a hard UI maximum limit, therefore the $max+$ boundary is impossible and omitted.}
\begin{table}[H]
\centering
\renewcommand{\arraystretch}{1.2}
\begin{tabularx}{\textwidth}{|c|c|l|X|}
\hline
\textbf{Boundary} & \textbf{Length} & \textbf{Target Variable Test Input} & \textbf{Expected Result} \\ \hline
$min-$ & 0 & \textit{(empty string)} & ❌ Error: Form submission fails \\ \hline
$min$ & 1 & \texttt{t} & ✅ Accepted: Event created \\ \hline
$min+$ & 2 & \texttt{te} & ✅ Accepted: Event created \\ \hline
\end{tabularx}
\end{table}

\vspace{0.3cm}
\noindent\textbf{Variable 2: Duration in Minutes} \\
\textit{Domain: $\text{Minutes} \ge 1$}
\begin{table}[H]
\centering
\renewcommand{\arraystretch}{1.2}
\begin{tabularx}{\textwidth}{|c|c|l|X|}
\hline
\textbf{Boundary} & \textbf{Value} & \textbf{Target Variable Test Input} & \textbf{Expected Result} \\ \hline
$min-$ & 0 & \texttt{0} & ❌ Error: Duration must be 1 or more \\ \hline
$min$ & 1 & \texttt{1} & ✅ Accepted: Event created \\ \hline
$min+$ & 2 & \texttt{2} & ✅ Accepted: Event created \\ \hline
$nom$ & 60 & \texttt{60} & ✅ Accepted: Event created \\ \hline
$max-$ & 9{,}999{,}998 & \texttt{9999998} & ✅ Accepted: Event created \\ \hline
$max$ & 9{,}999{,}999 & \texttt{9999999} & ✅ Accepted: Event created \\ \hline
$max+$ & 10{,}000{,}000 & \texttt{10000000} & ❌ Error: Duration must be 9{,}999{,}999 or less \\ \hline
\end{tabularx}
\end{table}

\subsubsection{Equivalence Class Partitioning (ECP)}
\textbf{Methodology Applied: Weak Robust ECP.}
The input domain is partitioned into mutually exclusive Valid and Invalid classes. One representative value is selected from each class. Similar to BVA, the Single-Fault Assumption is strictly enforced during execution.

\textbf{Variables Analyzed Summary:}
\begin{itemize}
    \item \textbf{V1: Duration Input Format} (3 Invalid Classes: DT4, DT5, DT6).
\end{itemize}

\vspace{0.3cm}
\noindent\textbf{Equivalence Classes Definition:}
\begin{table}[H]
\centering
\renewcommand{\arraystretch}{1.2}
\begin{tabularx}{\textwidth}{|c|l|X|l|X|}
\hline
\textbf{Var} & \textbf{ID} & \textbf{Equivalence Class} & \textbf{Representative} & \textbf{Expected Result} \\ \hline
\multirow{3}{*}{\textbf{Duration}}
& DT4 & Decimal (float) value & \texttt{1.5} & ❌ Error near Duration field \\ \cline{2-5}
& DT5 & Alphabetic characters & \texttt{abc} & ❌ Error near Duration field \\ \cline{2-5}
& DT6 & Negative value & \texttt{-10} & ❌ Error near Duration field \\ \hline
\end{tabularx}
\end{table}

\subsubsection{Decision Table Testing}
\textbf{Methodology Applied: Component Logic Mapping.} 
The duration settings interface relies on a set of mutually exclusive radio buttons. Selecting an option dynamically toggles the visibility and mandatory validation of corresponding fields.

\begin{table}[H]
\centering
\renewcommand{\arraystretch}{1.2}
\begin{tabularx}{\textwidth}{|>{\raggedright\arraybackslash}X|c|c|c|c|c|c|}
\hline
\textbf{Conditions} & \textbf{R1} & \textbf{R2} & \textbf{R3} & \textbf{R4} & \textbf{R5} & \textbf{R6} \\ \hline
C1: Duration mode & Without & Until & Minutes & Without & Until & Minutes \\ \hline
C2: Repeat this event & F & F & F & T & T & T \\ \hline
\hline
\textbf{Actions / System Output (UI Visibility)} & & & & & & \\ \hline
A1: All duration sub-fields hidden & \textbf{X} & & & \textbf{X} & & \\ \hline
A2: \textit{Until} Date/Time picker visible & & \textbf{X} & & & \textbf{X} & \\ \hline
A3: \textit{Minutes} number input visible & & & \textbf{X} & & & \textbf{X} \\ \hline
A4: Repeat count field visible & & & & \textbf{X} & \textbf{X} & \textbf{X} \\ \hline
\end{tabularx}
\end{table}

\subsubsection{Use-Case Testing}
\textbf{Methodology Applied: Path-based Scenario Generation.} 
Test scenarios are systematically derived by traversing distinct execution paths (Basic, Alternative, and Exception paths) identified in the Activity Diagram.

\vspace{0.2cm}
\noindent\textbf{Use Case Name:} User Creates Calendar Event \\
\textbf{Actor:} Authenticated User \\
\textbf{Pre-condition:} 
\begin{enumerate}
    \item Log in to \texttt{https://ihatetesting.moodlecloud.com/}\\ with username \texttt{phuc.nguyen0310@hcmut.edu.vn} and password \texttt{Huuphuc0310@}.
    \item Navigate to the \textbf{Calendar}.
    \item Click \textbf{New event}.
\end{enumerate}
\textbf{Post-condition:} A new calendar event is scheduled, followed by the user logging out of the system.

\vspace{0.3cm}
\noindent\textbf{Derived Test Scenarios:}
\begin{table}[H]
\centering
\renewcommand{\arraystretch}{1.2}
\begin{tabularx}{\textwidth}{|l|l|X|}
\hline
\textbf{Scenario ID} & \textbf{Path Type} & \textbf{Description \& Expected Result} \\ \hline
S1 & Basic Flow (Happy Path) & User enters valid title, date, selects "Duration in minutes" (60). \textbf{Result:} Event created successfully. \\ \hline
S2 & Alternative Flow & User leaves duration as "Without duration". \textbf{Result:} Event created successfully as an instantaneous event. \\ \hline
S3 & Alternative Flow & User selects "Until" and sets a valid future end date. \textbf{Result:} Event created successfully spanning multiple days. \\ \hline
S4 & Exception Flow & User leaves the Event Title field empty. \textbf{Result:} Form submission fails, error shown near Title. \\ \hline
S5 & Exception Flow & User sets "Duration in minutes" to 0 or negative. \textbf{Result:} Form submission fails, error shown near Duration. \\ \hline
\end{tabularx}
\end{table}

\begin{figure}[H]
    \centering
    % Chèn hình ảnh graphics.TC-001.png vào đây
    \includegraphics[width=0.8\textwidth]{graphics/TC-005.png}
    \caption{UML Activity Diagram outlining the execution paths for Calendar Event Creation}
    \label{fig:calendar-event-creation-activity}
\end{figure}

\subsubsection{Test Cases Summary (Excerpt)}
\textit{Note: This is an excerpt of key test cases. The complete set of test cases, detailing multi-variable data combinations, is documented in the attached \texttt{Testcases.xlsx} file.}

\begin{table}[H]
\centering
\renewcommand{\arraystretch}{1.2}
\begin{tabularx}{\textwidth}{|l|c|X|X|}
\hline
\textbf{TC ID} & \textbf{Method} & \textbf{Test Case Name / Coverage} & \textbf{Expected Result} \\ \hline
TC-005-002 & BVA & V1: Title Length (min-) & ❌ Error: Form submission fails \\ \hline
TC-005-005 & BVA & V2: Duration Minutes (min-) & ❌ Error: Form submission fails \\ \hline
TC-005-010 & BVA & V2: Duration Minutes (max+) & ❌ Error: Form submission fails \\ \hline
TC-005-016 & ECP & Duration (DT4: Float value) & ❌ Error: Form submission fails \\ \hline
TC-005-019 & DT  & R1: Without duration, No repeat & ✅ Accepted: Event created \\ \hline
TC-005-020 & DT  & R2: Until mode, No repeat & ✅ Accepted: Event created \\ \hline
\multicolumn{4}{|c|}{\cellcolor{gray!20}\textbf{Use-Case Scenarios Coverage}} \\ \hline
TC-005-025 & UC & Cover S1: Happy path (no duration) & ✅ Accepted: Event created \\ \hline
TC-005-026 & UC & Cover S2: With minutes duration & ✅ Accepted: Event created \\ \hline
TC-005-027 & UC & Cover S4: Empty title error & ❌ Error: Form submission fails \\ \hline
TC-005-028 & UC & Cover S5: Negative duration error & ❌ Error: Form submission fails \\ \hline
\end{tabularx}
\end{table}


% --- FEATURE 006 ---
\subsection{Feature 006: Teacher Sets Up a Quiz}

\subsubsection{Feature Description \& Usage Instruction}
The Quiz activity module allows teachers to design and build quizzes consisting of a large variety of question types. This feature contains critical input validations regarding grading rules, attempt limits, timing settings (open/close dates), and time limits.
\newline
\textbf{Usage Instruction:}
\begin{enumerate}
    \item Log in to \texttt{https://ihatetesting.moodlecloud.com/}\\ with username \texttt{phuc.nguyen0310@hcmut.edu.vn} and password \texttt{Huuphuc0310@}.
    \item Navigate to a specific course and toggle \textbf{Edit mode} on.
    \item Click \textbf{Add an activity or resource $>$ Quiz}.
    \item Fill in the Quiz name, Timing restrictions, and Grade configurations.
    \item Click \textbf{Save and return to course}.
\end{enumerate}

\subsubsection{Data Preparation \& Clean-up (Setup/Teardown)}
To ensure test independence and prevent data clutter during automated test execution, the following data lifecycle scripts must be executed for this feature:
\begin{itemize}
    \item \textbf{Setup Script:} Log in as a Teacher. Access the target test course. Ensure that no duplicate quiz names (e.g., \texttt{qn001\_QuizNom}) exist in the course sections. If they do, delete them to maintain a clean test environment.
    \item \textbf{Teardown Script:} Navigate back to the course page and ensure \textbf{Edit mode} is toggled on. Locate the newly created test quizzes. Click the 3-dots (More) menu next to each quiz, select \textbf{Delete}, and click \textbf{Yes} on the confirmation modal to permanently clear the data. Finally, \textbf{log out} of the system.
\end{itemize}

\subsubsection{Form Fields Catalog}
\begin{table}[H]
\centering
\renewcommand{\arraystretch}{1.3}
\begin{tabularx}{\textwidth}{|l|l|c|X|}
\hline
\textbf{Field} & \textbf{Type} & \textbf{Req.} & \textbf{Constraints / Business Rules} \\ \hline
Name & Text input & Yes & Cannot be empty. Maximum length enforced by UI. \\ \hline
Open the quiz & Checkbox / Date & No & If enabled, must be a valid calendar date and time. \\ \hline
Close the quiz & Checkbox / Date & No & If enabled, must be strictly ON or AFTER the "Open the quiz" date. \\ \hline
Time limit & Checkbox / Number & No & If enabled, value must be $\ge 1$ (minutes, hours, days, etc.). \\ \hline
Grade to pass & Number input & No & Must be a valid number $\le$ Maximum grade (Default maximum is 10.00). \\ \hline
\end{tabularx}
\end{table}

\subsubsection{Boundary Value Analysis (BVA)}
\textbf{Methodology Applied: Robust Boundary Value Analysis (Robust BVA).} 
Based on the single-fault assumption theory, for a variable bounded by $[min, max]$, Robust BVA generates exactly 7 test values. \textbf{Critical Execution Note:} When the target variable is tested at a boundary, all other variables in the form are explicitly held at their valid nominal ($nom$) values to prevent fault masking.

\textbf{Variables Analyzed Summary:}
\begin{itemize}
    \item \textbf{V1: Grade to Pass} (3 values: $max-, max, max+$ --- lower boundary unconstrained by UI; $min-$/$min$/$min+$/$nom$ omitted).
    \item \textbf{V2: Time Limit} (7 values: $min-, min, min+, nom, max-, max, max+$ --- bounding domain minutes).
    \item \textbf{V3: Close Date Chronology} (6 values: $min-, min, min+, max-, max, max+$ relative to Open Date).
\end{itemize}

\vspace{0.3cm}
\noindent\textbf{Variable 1: Grade to Pass} \\
\textit{Domain: $0 \le \text{Grade} \le 10.00$}
\begin{table}[H]
\centering
\renewcommand{\arraystretch}{1.2}
\begin{tabularx}{\textwidth}{|c|c|l|X|}
\hline
\textbf{Boundary} & \textbf{Value} & \textbf{Target Variable Test Input} & \textbf{Expected Result} \\ \hline
$max-$ & 9.99 & \texttt{9.99} & ✅ Accepted: Quiz created successfully \\ \hline
$max$ & 10.00 & \texttt{10} & ✅ Accepted: Quiz created successfully \\ \hline
$max+$ & 10.01 & \texttt{10.01} & ❌ Error: Grade to pass cannot be greater than maximum possible grade \\ \hline
\end{tabularx}
\end{table}

\vspace{0.3cm}
\noindent\textbf{Variable 2: Time Limit} \\
\textit{Domain: $\text{Minutes} \ge 0$ (0 = unlimited); negative values are rejected.}
\begin{table}[H]
\centering
\renewcommand{\arraystretch}{1.2}
\begin{tabularx}{\textwidth}{|c|c|l|X|}
\hline
\textbf{Boundary} & \textbf{Value} & \textbf{Target Variable Test Input} & \textbf{Expected Result} \\ \hline
$min-$ & -1 & \texttt{-1} & ❌ Error: Time limit must be 0 or more \\ \hline
$min$ & 0 & \texttt{0} & ✅ Accepted: Quiz created successfully \\ \hline
$min+$ & 1 & \texttt{1} & ✅ Accepted: Quiz created successfully \\ \hline
$nom$ & 30 & \texttt{30} & ✅ Accepted: Quiz created successfully \\ \hline
$max-$ & 998 & \texttt{998} & ✅ Accepted: Quiz created successfully \\ \hline
$max$ & 999 & \texttt{999} & ✅ Accepted: Quiz created successfully \\ \hline
$max+$ & 1000 & \texttt{1000} & ✅ Accepted: Quiz created successfully \\ \hline
\end{tabularx}
\end{table}

\vspace{0.3cm}
\noindent\textbf{Variable 3: Close Date Chronology (Relative to Open Date)}
\begin{table}[H]
\centering
\renewcommand{\arraystretch}{1.2}
\begin{tabularx}{\textwidth}{|c|l|X|}
\hline
\textbf{Boundary} & \textbf{Target Variable Test Input} & \textbf{Expected Result} \\ \hline
$min-$ & \texttt{Open Date - 1 day} & ❌ Error: Close date must be after open date \\ \hline
$min$ & \texttt{Exact Open Date} & ✅ Accepted: Quiz created successfully \\ \hline
$min+$ & \texttt{Open Date + 1 day} & ✅ Accepted: Quiz created successfully \\ \hline
$max-$ & \texttt{Open Date + 10 years} & ✅ Accepted: Quiz created successfully \\ \hline
$max$ & \texttt{Open Date + 11 years} & ✅ Accepted: Quiz created successfully \\ \hline
$max+$ & \texttt{Open Date + 12 years} & ✅ Accepted: Quiz created successfully \\ \hline
\end{tabularx}
\end{table}

\subsubsection{Equivalence Class Partitioning (ECP)}
\textbf{Methodology Applied: Weak Robust ECP.}
The input domain is partitioned into mutually exclusive Valid and Invalid classes. One representative value is selected from each class. Similar to BVA, the Single-Fault Assumption is strictly enforced during execution.

\vspace{0.3cm}
\noindent\textbf{Equivalence Classes Definition:}
\begin{table}[H]
\centering
\renewcommand{\arraystretch}{1.2}
\begin{tabularx}{\textwidth}{|c|l|X|l|X|}
\hline
\textbf{Var} & \textbf{ID} & \textbf{Equivalence Class} & \textbf{Representative} & \textbf{Expected Result} \\ \hline
\textbf{Quiz Name} & QN3 & Empty quiz name & \textit{(empty)} & ❌ Error: Quiz creation fails \\ \hline
\end{tabularx}
\end{table}

\subsubsection{Decision Table Testing}
\textbf{Methodology Applied: Combinatorial Logic Modeling.} 
The Quiz timing settings feature three independent functional checkboxes ("Enable Open date", "Enable Close date", "Enable Time limit"). Based on the project report specifications, 5 distinct rules cover the primary operational modes of the Quiz timing configuration.

\begin{table}[H]
\centering
\renewcommand{\arraystretch}{1.2}
\begin{tabularx}{\textwidth}{|>{\raggedright\arraybackslash}X|c|c|c|c|c|}
\hline
\textbf{Conditions} & \textbf{R1} & \textbf{R2} & \textbf{R3} & \textbf{R4} & \textbf{R5} \\ \hline
C1: Open date enabled & T & T & T & F & F \\ \hline
C2: Close date enabled & T & T & F & T & F \\ \hline
C3: Time limit enabled & T & F & T & T & F \\ \hline
\hline
\textbf{Actions / System Output} & & & & & \\ \hline
A1: Allow access only within dates \& limit time & \textbf{X} & & & & \\ \hline
A2: Allow access only within dates (no timer) & & \textbf{X} & & & \\ \hline
A3: Allow access strictly after open date \& limit time & & & \textbf{X} & & \\ \hline
A4: Allow access strictly before close date \& limit time & & & & \textbf{X} & \\ \hline
A5: Quiz always open without restrictions & & & & & \textbf{X} \\ \hline
\end{tabularx}
\end{table}

\subsubsection{Use-Case Testing}
\textbf{Methodology Applied: Path-based Scenario Generation.} 
Test scenarios are systematically derived by traversing distinct execution paths (Basic, Alternative, and Exception paths) identified in the Activity Diagram.

\vspace{0.2cm}
\noindent\textbf{Use Case Name:} Teacher Sets Up a Quiz \\
\textbf{Actor:} Teacher \\
\textbf{Pre-condition:} 
\begin{enumerate}
    \item Log in to \texttt{https://ihatetesting.moodlecloud.com/}\\ with username \texttt{phuc.nguyen0310@hcmut.edu.vn} and password \texttt{Huuphuc0310@}.
    \item Navigate to the target course and toggle \textbf{Edit mode} on.
    \item Click \textbf{Add an activity or resource $>$ Quiz}.
\end{enumerate}
\textbf{Post-condition:} A new quiz activity is created and visible in the course sections, followed by the user logging out of the system.

\vspace{0.3cm}
\noindent\textbf{Derived Test Scenarios:}
\begin{table}[H]
\centering
\renewcommand{\arraystretch}{1.2}
\begin{tabularx}{\textwidth}{|l|l|X|}
\hline
\textbf{Scenario ID} & \textbf{Path Type} & \textbf{Description \& Expected Result} \\ \hline
S1 & Basic Flow (Happy Path) & Teacher configures all timing explicitly (R1) with valid Grade to Pass. \textbf{Result:} Quiz created successfully. \\ \hline
S2 & Alternative Flow & Unlimited + no time: All timing limits disabled (R5). \textbf{Result:} Quiz created successfully (always open). \\ \hline
S3 & Exception Flow & Date conflict: Close date is set earlier than Open date. \textbf{Result:} Form submission fails, error shown. \\ \hline
S4 & Exception Flow & Grade exceeds max: Grade to pass entered is greater than the Maximum grade (e.g., 100 on a 10 pt scale). \textbf{Result:} Form submission fails, error shown. \\ \hline
S5 & Exception Flow & Empty required field: Quiz Name left blank. \textbf{Result:} Form submission fails, error shown near Name. \\ \hline
\end{tabularx}
\end{table}

\begin{figure}[H]
    \centering
    % Chèn hình ảnh graphics.TC-001.png vào đây
    \includegraphics[width=0.8\textwidth]{graphics/TC-006.png}
    \caption{UML Activity Diagram outlining the execution paths for Quiz Creation}
    \label{fig:quiz-creation-activity}
\end{figure}

\subsubsection{Test Cases Summary (Excerpt)}
\textit{Note: This is an excerpt of key test cases. The complete set of test cases, detailing multi-variable data combinations, is documented in the attached \texttt{Testcases.xlsx} file.}

\begin{table}[H]
\centering
\renewcommand{\arraystretch}{1.2}
\begin{tabularx}{\textwidth}{|l|c|X|X|}
\hline
\textbf{TC ID} & \textbf{Method} & \textbf{Test Case Name / Coverage} & \textbf{Expected Result} \\ \hline
TC-006-001 & BVA & V1: All (nom) & ✅ Accepted: Quiz created successfully \\ \hline
TC-006-002 & BVA & V1: Grade to Pass (max-) & ✅ Accepted: Quiz created successfully \\ \hline
TC-006-011 & BVA & V3: Close Date (min-) & ❌ Error: Form submission fails \\ \hline
TC-006-017 & ECP & Quiz Name (QN3: Empty) & ❌ Error: Form submission fails \\ \hline
\multicolumn{4}{|c|}{\cellcolor{gray!20}\textbf{Decision Table \& Use-Case Scenarios Coverage}} \\ \hline
TC-006-018 & DT  & R1: All timing enabled & ✅ Accepted: Quiz created successfully \\ \hline
TC-006-019 & DT  & R2: Open+Close only & ✅ Accepted: Quiz created successfully \\ \hline
TC-006-022 & DT  & R5: All timing disabled & ✅ Accepted: Quiz created successfully \\ \hline
TC-006-023 & UC  & Cover S1: Happy path S1 & ✅ Accepted: Quiz created successfully \\ \hline
TC-006-024 & UC  & Cover S2: Unlimited + no time S2 & ✅ Accepted: Quiz created successfully \\ \hline
TC-006-025 & UC  & Cover S3: Date conflict S3 & ❌ Error: Form submission fails \\ \hline
TC-006-026 & UC  & Cover S4: Grade exceeds max S4 & ❌ Error: Form submission fails \\ \hline
\end{tabularx}
\end{table}

\newpage
% % ==========================================
% % 4. TEST AUTOMATION & EXECUTION
% % ==========================================
% \section{Test Automation \& Execution}

% \subsection{Data Preparation and Clean-up Scripts (Setup/Teardown)}
% To ensure test independence and system integrity, automated scripts were utilized before and after executing the main test suites in Katalon Recorder:
% \begin{itemize}
%     \item \textbf{Setup Script (Pre-condition):} Executes before the main test suite to prepare the environment. For example, a script logs in as an admin to pre-create a prerequisite Category or a baseline user account required for testing Feature 001.
%     \item \textbf{Teardown Script (Post-condition):} Executes after the test suite finishes to clean up the data. It automates the deletion of the newly generated test courses or dummy users (e.g., \texttt{testuser01}) to restore Moodle to its original state, preventing data clutter for subsequent test runs.
% \end{itemize}

% ==========================================
% 5. CONCLUSION
% ==========================================
\section{Conclusion}
This project successfully applied four major black-box testing methodologies to the Moodle system. By integrating Boundary Value Analysis, Equivalence Class Partitioning, Decision Tables, and Use-case Testing, the team achieved high test coverage. 

Through this testing campaign, our team gained profound practical insights into software verification. We successfully translated theoretical test designs into automated test suites using Katalon Recorder, significantly improving execution efficiency. The challenge of migrating test environments due to unexpected server errors also taught us the critical importance of adaptability in real-world testing scenarios. Overall, this project solidified our understanding of quality assurance workflows and enhanced our technical problem-solving capabilities.

\end{document}